\chapter{Introduction to Medical Testing}

Medical testing is the bridge between what a clinician observes and what is actually happening inside the body. A patient's history and physical examination generate hypotheses; laboratory and imaging tests test those hypotheses with objective data. This chapter provides an overview of what medical tests measure, how they inform clinical decision-making, and how this book is structured to serve as a practical reference.

\clearpage
\section{The Purpose of Medical Testing}

Every medical test is designed to answer a specific clinical question. These questions fall into several broad categories:

\begin{itemize}
    \item \textbf{Detection:} Does the patient have a particular condition? Examples include troponin for myocardial infarction, PCR for SARS-CoV-2, or mammography for breast cancer.
    \item \textbf{Severity assessment:} How advanced is the disease? Examples include HbA1c for glycemic control, eGFR for kidney function, or tumor staging via CT and PET.
    \item \textbf{Monitoring:} Is the condition improving, worsening, or stable? Examples include serial CBC during chemotherapy, INR during anticoagulation, or viral load during antiretroviral therapy.
    \item \textbf{Screening:} Does the patient have risk factors or early-stage disease before symptoms appear? Examples include lipid panels for cardiovascular risk, Pap smears for cervical cancer, or newborn metabolic screens.
    \item \textbf{Prognosis:} What is the likely outcome? Examples include BNP in heart failure, PSA velocity in prostate cancer, or cytogenetic markers in leukemia.
\end{itemize}

A single test may serve multiple purposes depending on the clinical context. The same HbA1c used for screening (detecting undiagnosed diabetes) is also used for monitoring (assessing glycemic control over time).

\clearpage
\section{What Medical Tests Measure}

Despite the diversity of available tests, they all measure a limited set of fundamental biological signals. Understanding these categories helps in interpreting any test result:

\subsection{Chemical Constituents}

The largest category of laboratory tests measures the concentration of chemical substances in blood, urine, or other body fluids:

\begin{itemize}
    \item \textbf{Electrolytes and metabolites:} Sodium, potassium, chloride, bicarbonate, glucose, creatinine, urea, lactate. These reflect homeostasis, organ function, and metabolic state. Abnormalities can indicate dehydration, renal failure, diabetic ketoacidosis, or respiratory disorders.
    \item \textbf{Enzymes:} ALT, AST, ALP, GGT, CK, LDH, amylase, lipase. These are measured as indicators of cellular injury. When cells die or become leaky, intracellular enzymes spill into the bloodstream. The pattern of enzyme elevation identifies the damaged organ.
    \item \textbf{Hormones:} TSH, free T4, cortisol, PTH, insulin, estradiol, testosterone. Hormone levels reflect endocrine gland function and feedback regulation. Abnormalities can indicate primary gland dysfunction or secondary disruption of regulatory axes.
    \item \textbf{Proteins:} Albumin, globulins, prealbumin, transferrin, ferritin, CRP, complement. These serve diverse functions including transport, immunity, inflammation, and nutrition assessment.
    \item \textbf{Lipids:} Total cholesterol, LDL, HDL, triglycerides. These are measured primarily for cardiovascular risk assessment.
    \item \textbf{Drugs and toxins:} Therapeutic drug levels (digoxin, lithium, vancomycin), drugs of abuse, toxic alcohols, heavy metals. These confirm exposure, guide dosing, and identify poisoning.
\end{itemize}

\subsection{Cellular Elements}

Blood and body fluids contain cells whose number, morphology, and function provide diagnostic information:

\begin{itemize}
    \item \textbf{Complete blood count:} RBC, WBC, hemoglobin, hematocrit, platelet count, differential. This is the most frequently ordered test in medicine. It detects anemia, infection, inflammation, bleeding disorders, and hematologic malignancies.
    \item \textbf{Peripheral smear:} Visual examination of blood cells under a microscope. Identifies abnormal cell shapes, immature cells, parasites, and confirms automated count abnormalities.
    \item \textbf{Body fluid cell counts:} Cells in CSF, pleural fluid, ascitic fluid, synovial fluid. Differentiates transudate from exudate, identifies infection, and detects malignant cells.
    \item \textbf{Bone marrow aspirate/biopsy:} Evaluates hematopoietic tissue for leukemia, lymphoma, myeloma, aplastic anemia, and metastatic disease.
\end{itemize}

\subsection{Immunologic Markers}

The immune system produces molecules that can be measured to diagnose infectious, autoimmune, allergic, and neoplastic conditions:

\begin{itemize}
    \item \textbf{Antibodies:} Pathogen-specific (HBsAg, anti-HCV, anti-HIV), autoantibodies (ANA, RF, anti-CCP, anti-dsDNA), allergenspecific IgE. Serologic testing detects past or current infection and identifies autoimmune disease.
    \item \textbf{Antigens:} Pathogen proteins (p24 for HIV, NS1 for dengue), tumor antigens (PSA, CA-125, CEA). Direct detection of microbial or tumor products.
    \item \textbf{Acute phase reactants:} CRP, ESR, procalcitonin, ferritin. Nonspecific markers of inflammation that rise in response to tissue injury, infection, or autoimmune flares.
    \item \textbf{Cytokines and chemokines:} IL-6, TNF-alpha, IFN-gamma. Measured in specialized contexts for immune monitoring and research.
\end{itemize}

\subsection{Genetic Material}

Molecular and cytogenetic testing examines DNA and RNA to identify inherited mutations, somatic mutations, chromosomal abnormalities, and microbial genetic material:

\begin{itemize}
    \item \textbf{Chromosomal analysis:} Karyotyping identifies structural abnormalities in chromosomes (translocations, deletions, duplications) in constitutional disorders and malignancies.
    \item \textbf{Gene sequencing:} Targeted panels, exome sequencing, or genome sequencing identify specific mutations causing genetic disorders or driving cancer.
    \item \textbf{Expression analysis:} RNA-based assays measure gene expression levels to classify tumors and predict treatment response.
    \item \textbf{Microbial nucleic acid detection:} PCR and related methods detect pathogen DNA or RNA directly, enabling diagnosis of infections that are difficult to culture.
\end{itemize}

\subsection{Physiologic Function}

Some tests assess how well an organ or system performs its function rather than measuring a single molecule:

\begin{itemize}
    \item \textbf{Renal function:} eGFR, creatinine clearance, urine protein-to-creatinine ratio. These estimate the kidney's ability to filter blood.
    \item \textbf{Hepatic function:} Bilirubin, albumin, PT/INR, ammonia. These measure the liver's synthetic and excretory capacity.
    \item \textbf{Cardiac function:} Ejection fraction (echocardiography), BNP/NT-proBNP, exercise stress testing, coronary angiography. These evaluate cardiac mechanics, perfusion, and hemodynamics.
    \item \textbf{Pulmonary function:} Spirometry, DLCO, ABG. These measure lung mechanics, gas exchange, and oxygenation.
\end{itemize}

\subsection{Anatomic Structure}

Imaging modalities provide structural and sometimes functional information about organs and tissues:

\begin{itemize}
    \item \textbf{X-ray:} Density-based projection imaging. Primary tool for bone fractures, pneumonia, pneumothorax, and bowel obstruction.
    \item \textbf{CT:} Cross-sectional imaging with excellent spatial resolution. Used for trauma, cancer staging, pulmonary embolism, and intra-abdominal pathology.
    \item \textbf{MRI:} High-contrast soft tissue imaging. Essential for brain, spine, joint, and pelvic pathology.
    \item \textbf{Ultrasound:} Real-time imaging without ionizing radiation. Used for obstetrics, cardiac function, abdominal organs, and vascular assessment.
    \item \textbf{Nuclear medicine (PET/SPECT):} Functional imaging based on metabolic activity. Critical for oncology, infection imaging, and myocardial perfusion.
\end{itemize}

\clearpage
\section{From Measurement to Clinical Understanding}

A test result becomes clinically useful only when interpreted in context. Several principles govern how measurements translate into understanding of disease:

\subsection{Reference Ranges and Variability}

A test result is compared to a reference range derived from a healthy population. However:
\begin{itemize}
    \item Reference ranges vary by age, sex, ethnicity, and laboratory methodology.
    \item A result within the reference range does not rule out disease, and a result outside it does not confirm disease.
    \item Serial results in the same individual are often more informative than a single measurement, because each person has their own baseline.
\end{itemize}

\subsection{Pattern Recognition}

Few diseases are diagnosed by a single test. Most are recognized by characteristic patterns across multiple tests. For example:
\begin{itemize}
    \item \textbf{Hepatocellular injury:} Elevated ALT and AST with normal or mildly elevated ALP and GGT.
    \item \textbf{Cholestatic injury:} Elevated ALP and GGT with mild ALT/AST elevation.
    \item \textbf{Iron deficiency anemia:} Low hemoglobin, low MCV, low ferritin, low iron, high TIBC.
    \item \textbf{Anemia of chronic disease:} Low hemoglobin, low MCV, normal/high ferritin, low iron, low TIBC.
\end{itemize}
The pattern, not any single value, drives the diagnosis.

\subsection{Sensitivity, Specificity, and Predictive Value}

Every test has performance characteristics that determine how much trust to place in its result:
\begin{itemize}
    \item \textbf{Sensitivity:} The proportion of true positives correctly identified. A highly sensitive test rarely misses disease.
    \item \textbf{Specificity:} The proportion of true negatives correctly identified. A highly specific test rarely gives false positives.
    \item \textbf{Predictive values:} Depend on disease prevalence. Even a highly accurate test yields many false positives when the disease is rare (screening), and many false negatives when the disease is common.
\end{itemize}

\subsection{Trends Over Time}

A single measurement is a snapshot. Serial measurements reveal trajectories that are often more meaningful:
\begin{itemize}
    \item Rising troponin over 6 hours confirms myocardial infarction.
    \item Falling eGFR over months indicates progressive kidney disease.
    \item Rising PSA over years may indicate progressing prostate cancer.
    \item Falling viral load confirms response to antiretroviral therapy.
\end{itemize}

\clearpage
\section{How This Book Is Organized}

Each chapter in this book covers a major domain of medical testing. Within each chapter, individual tests are presented in a standardized format:

\begin{itemize}
    \item \textbf{Test name and synonyms}
    \item \textbf{Specimen type and collection requirements}
    \item \textbf{Analytical method}
    \item \textbf{Turnaround time}
    \item \textbf{Reference range}
    \item \textbf{Clinical interpretation} --- what the test measures, why it is ordered, and how results are interpreted
    \item \textbf{Factors affecting results} --- interferences, medications, physiologic states
    \item \textbf{Follow-up testing} --- what additional tests may be indicated based on results
\end{itemize}

The appendices provide reference ranges, critical values, test interference information, billing codes, diagnostic algorithms, and quality control procedures for rapid reference.

\subsection{A Note on Cost}

In the United States, the cost of laboratory tests varies enormously depending on where they are performed. The same CBC that costs \$10 at a hospital lab can cost \$150 at a freestanding commercial lab. If cost is a concern, always ask:

\begin{itemize}
    \item Is this test covered by my insurance?
    \item Can it be performed at a lower-cost facility?
    \item Are there alternative tests that provide the same information for less?
\end{itemize}

Do not let financial barriers prevent you from getting the care you need. Most hospitals have financial assistance programs. Many laboratories offer cash-pay discounts. Ask before you consent.
